\section{SPDA e proteção contra descargas atmosféricas}

\begin{frame}{NBR 5419: atualização confirmada em 2026}
\small
\begin{columns}[T,onlytextwidth]
\column{0.52\textwidth}
\begin{caixaNorma}[title=ABNT NBR 5419-1:2026]
A edição 2026 da \textbf{Parte 1} está publicada e estabelece os requisitos para a determinação da necessidade de proteção contra descargas atmosféricas.
\end{caixaNorma}
\begin{itemize}
\item não assumir automaticamente que todas as partes da série mudaram de edição na mesma data;
\item no projeto, registrar a edição efetivamente adotada de cada parte aplicável;
\item verificar requisitos do órgão licenciador, Corpo de Bombeiros e contratante.
\end{itemize}
\column{0.46\textwidth}
\begin{caixaAt}[title=Regra para este material]
A referência ``NBR 5419:2015'' não é mais usada de forma genérica. Quando a Parte 1 for citada, o material indica explicitamente \textbf{ABNT NBR 5419-1:2026}; para as demais partes, exige-se conferência da edição vigente antes do projeto executivo.
\end{caixaAt}
\end{columns}
\end{frame}

\begin{frame}{PDA não é apenas o captor no telhado}
\small
\centering
\begin{tikzpicture}[scale=0.82,every node/.style={font=\scriptsize}]
\tikzset{b/.style={draw=corPrim,fill=corDef!7,rounded corners,minimum width=3.0cm,minimum height=0.7cm,align=center}}
\node[b] (r) at (0,2.6) {Avaliação da necessidade\\e análise de risco};
\node[b] (s) at (-3.4,1.1) {SPDA externo\\captação/descida/terra};
\node[b] (e) at (0,1.1) {Equipotencialização\\para descargas};
\node[b] (m) at (3.4,1.1) {MPS\\DPS, roteamento, blindagem};
\node[b] (i) at (0,-0.5) {Sistemas internos\\energia, dados e automação};
\draw[-{Stealth},thick,corPrim] (r)--(s);\draw[-{Stealth},thick,corPrim] (r)--(e);\draw[-{Stealth},thick,corPrim] (r)--(m);
\draw[-{Stealth},thick,corNorma] (s)--(i);\draw[-{Stealth},thick,corNorma] (e)--(i);\draw[-{Stealth},thick,corNorma] (m)--(i);
\end{tikzpicture}
\vspace{2mm}
\begin{caixaObs}
A interface entre SPDA, equipotencialização e DPS precisa ser projetada conjuntamente. Separar esses projetos em ``ilhas'' cria lacunas de proteção.
\end{caixaObs}
\end{frame}

\begin{frame}{Zonas de proteção e roteamento}
\small
\begin{columns}[T,onlytextwidth]
\column{0.49\textwidth}
\begin{itemize}
\item A energia do surto deve ser reduzida progressivamente entre zonas.
\item Entradas metálicas e cabos são pontos críticos de acoplamento.
\item Roteamento, blindagem, equipotencialização e DPS trabalham em conjunto.
\item Cabos sensíveis não devem percorrer rotas inadequadas junto a condutores de descida sem análise.
\end{itemize}
\column{0.49\textwidth}
\centering
\begin{tikzpicture}[scale=0.82,every node/.style={font=\scriptsize}]
\draw[thick,corAt] (0,0) rectangle (5.3,3.4);
\draw[thick,corAccent] (0.7,0.6) rectangle (4.6,2.8);
\draw[thick,corNorma] (1.5,1.2) rectangle (3.8,2.2);
\node[corAt] at (4.65,3.1) {ZPR externa};
\node[corAccent] at (3.9,2.55) {zona interna};
\node[corNorma] at (2.65,1.7) {equipamento sensível};
\draw[-{Stealth},thick,corPrim] (-0.6,1.7)--(1.4,1.7) node[midway,above]{linha};
\node at (0.7,1.35) {DPS};
\end{tikzpicture}
\end{columns}
\end{frame}

%=====================================================================